% Источники: exam/materials/преподаватель/2020/20-04-2020-Формализация_математической_модели___пример_к_лекции_9_.pdf; exam/materials/моделирование.pdf, разделы 14 и 48.
\section{Методика построения аналитической статической модели.}

\textbf{Статическая модель} описывает состояние системы в фиксированный момент времени. Применяется при сравнении «срезов» динамической модели в разные моменты, а также при сложности динамических моделей или для наглядности представления. Примеры: принципиальные схемы устройств, цифровые карты местности, модели поверхностей.

Статическая модель может быть иллюстративной частью проекта, использующего физическую или динамическую математическую модель.

\subsection*{Методика построения (7 шагов)}

\begin{enumerate}
  \item Рисуется \textbf{графическая схема}, иллюстрирующая состояние системы в некоторый момент времени в заданной системе координат.

  \item По схеме \textbf{уточняются входные данные} (например: матрица высот), погрешность представления измеряемой координаты; выделяются промежуточные переменные для оценки выходных данных; на схеме выделяется зависимость исследуемого показателя от остальных характеристик.

  \item \textbf{Определяются физико-химические воздействия} на объект, переводящие систему из предыдущего в текущее состояние.

  \item \textbf{Оценивается возможность формализации модели и её упрощения.} Необходимо обосновать возможность описания и целесообразность учёта выбранных физико-химических воздействий на исследуемый объект. В отличие от динамических моделей, в статических большее значение имеет учёт \textbf{химических} воздействий.

  \item \textbf{Формулируется общее уравнение} на языке математики, описывающее состояние объекта в заданный период времени, с учётом выделенных физико-химических воздействий.

  \item \textbf{Конкретизируется описание} состояния объекта: однозначно выделяются входные, выходные, промежуточные переменные и параметры модели.

  \item При формализации модели должно выполняться \textbf{условие единственности решения}; в силу чего ряд входных данных может быть доопределён либо исключён.
\end{enumerate}

Статическое моделирование описывает процесс дискретно, не давая полной картины его развития, но демонстрирует качественное (а не количественное) решение при низкой ресурсозатратности.

\clearpage
